% ===============================================
% MATH 3053: Abstract Algebra I         Fall 2017
% hw_revised_cn.tex
% 中文适配版作业修订提交模板
% ===============================================
%         请仔细阅读以下说明！！！
% ===============================================
% 生成PDF提交时，将文件名改为 hwX-姓名.pdf
% X替换为作业编号，姓名替换为你的姓氏
% ===============================================


% -----------------------------------------------
% 以下序言部分可忽略，直接跳至 "START HERE" 区域
% -----------------------------------------------

\documentclass{article}
% 页面设置与基础宏包
\usepackage[margin=1in]{geometry}  % 页边距1英寸
\usepackage{amsmath,amsthm,amssymb} % 数学公式、定理环境、符号
\usepackage{ctex}   % 中文支持核心宏包（Overleaf预装）
\usepackage{pifont}                 
\renewcommand{\qedsymbol}{}
\usepackage{booktabs}
\usepackage{graphicx} 
\usepackage{array}
% 常用数学集合符号定义
\newcommand{\R}{\mathbb{R}}  % 实数集
\newcommand{\Z}{\mathbb{Z}}  % 整数集
\newcommand{\N}{\mathbb{N}}  % 自然数集
\newcommand{\Q}{\mathbb{Q}}  % 有理数集
\newcommand{\C}{\mathbb{C}}  % 复数集
%代码框
% 代码框与语法高亮所需宏包
\usepackage{listings}    % 核心代码框宏包
\usepackage{xcolor}      % 支持颜色，用于语法高亮
\usepackage{geometry}    % 可选，调整页面边距，避免代码溢出
\usepackage{float} 
% 配置Python语法高亮样式
\lstset{
    language=Python,             % 指定代码语言为Python
    basicstyle=\ttfamily\small,  % 基础字体：等宽字体+小字号
    keywordstyle=\color{blue},   % 关键字颜色（如if/for/def）
    commentstyle=\color{gray},   % 注释颜色
    stringstyle=\color{green},   % 字符串颜色
    numbers=left,                % 行号显示在左侧
    numberstyle=\tiny\color{gray},% 行号样式：小字号+灰色
    frame=tb,                    % 代码框：上下边框（t=top, b=bottom）
    tabsize=4,                   % 制表符宽度（匹配Python标准）
    showstringspaces=false,      % 不显示字符串中的空格标记
    breaklines=true,             % 自动换行（避免代码溢出）
    breakatwhitespace=false,     % 可在非空格处换行（保证代码完整性）
    escapeinside={\%*}{*)},      % 可选，支持在代码中插入LaTeX命令
    captionpos=t,                % 标题位置在代码框顶部
    abovecaptionskip=0.5em,      % 标题与代码框的间距
    belowskip=1em                % 代码框与下文的间距
}
% 定理、习题等环境定义（支持中英文标注）
\newenvironment{theorem}[2][Theorem]{\begin{trivlist}
\item[\hskip \labelsep {\bfseries #1}\hskip \labelsep {\bfseries #2.}]}{\end{trivlist}}
\newenvironment{lemma}[2][Lemma]{\begin{trivlist}
\item[\hskip \labelsep {\bfseries #1}\hskip \labelsep {\bfseries #2.}]}{\end{trivlist}}
\newenvironment{exercise}[2][Exercise]{\begin{trivlist}
\item[\hskip \labelsep {\bfseries #1}\hskip \labelsep {\bfseries #2.}]}{\end{trivlist}}
% 中文"问题"环境（与原problem并存，可按需选用）
\newenvironment{problem}[2][Problem]{\begin{trivlist}
\item[\hskip \labelsep {\bfseries #1}\hskip \labelsep {\bfseries #2.}]}{\end{trivlist}}
\newenvironment{problem_cn}[2][问题]{\begin{trivlist}
\item[\hskip \labelsep {\bfseries #1}\hskip \labelsep {\bfseries #2.}]}{\end{trivlist}}
\newenvironment{question}[2][Question]{\begin{trivlist}
\item[\hskip \labelsep {\bfseries #1}\hskip \labelsep {\bfseries #2.}]}{\end{trivlist}}
\newenvironment{corollary}[2][Corollary]{\begin{trivlist}
\item[\hskip \labelsep {\bfseries #1}\hskip \labelsep {\bfseries #2.}]}{\end{trivlist}}

% 解答环境（支持中英文标注）
\newenvironment{solution}{\begin{proof}[Solution]}{\end{proof}}
\newenvironment{solution_cn}{%
  \begin{proof}[解答]%
  \kaishu  % 切换中文为楷体（ctex 提供的命令）
}{%
  \end{proof}%
}

\begin{document}

% ------------------------------------------ %
%                 从这里开始编辑                 %
% ------------------------------------------ %

% 标题与作者信息（中文适配）
\title{《优化方法基础》课程作业\\等式不等式约束优化问题}  % X替换为作业编号
\author{姓名：林圣翔 \quad 班级：计试2201 \quad 学号：2223312202}  % 替换括号内容

\maketitle

\begin{problem_cn}{1}
总结以下不等式约束凸优化问题的求解方法及相互联系
\[
\begin{aligned}
&\min && f_0(x) \\
&\text{s.t.} && f_i(x) \leq 0,\, i = 1,2,\cdots, m \\
&&& Ax = b
\end{aligned}
\]

\begin{itemize}
    \item 定义域 \( D = \bigcap_{i=0}^m \text{dom}\, f_i \);
    \item \( A \in \mathbb{R}^{p \times n} \)，\(\text{rank}(A) = p < n \);
    \item 函数 \( f_i: \mathbb{R}^n \to \mathbb{R} \, (i = 0,\cdots, m) \) 为二次连续可微凸函;
    \item \( p^* = \inf \{ f_0(x) : Ax = b,\, f_i(x) \leq 0,\, i = 1,\cdots, m \} = f(x^*) \)
\end{itemize}
\end{problem_cn}
\begin{solution_cn} 
\quad \\1.障碍方法

障碍方法将不等式约束转化为障碍函数，从而将原问题转化为一系列等式约束优化问题。\\
1.1.障碍函数

使用对数障碍函数近似示性函数：$ \phi(x) = -\sum_{i=1}^m \log\bigl(-f_i(x)\bigr)$ 

近似示性函数为 \(\hat{I}(u) = -\dfrac{1}{t} \log\bigl(-u\bigr)\)。\\
1.2.近似问题

对于参数 \(t > 0\)，求解$\min\ t f_0(x) + \phi(x) \quad \text{s.t.} \quad A x = b $ 。

该问题所得的解 \(x^*(t)\) 称为中心点。\\
1.3.中心路径

中心点 \(x^*(t)\) 满足严格可行性（\(A x^*(t) = b,\ f_i(x^*(t)) < 0\)）和最优性条件：  
  \[
  t \nabla f_0\bigl(x^*(t)\bigr) + \sum_{i=1}^m \frac{1}{-f_i\bigl(x^*(t)\bigr)} \nabla f_i\bigl(x^*(t)\bigr) + A^\top v = 0
  \]  

从中心点可导出对偶变量 \(\lambda_i^*(t) = \frac{1}{-t f_i\bigl(x^*(t)\bigr)} > 0\)，且对偶间隙满足：  
  \[
  f_0\bigl(x^*(t)\bigr) - p^* \leq \frac{m}{t}
  \]  
1.4.算法流程

1) 初始化：找到严格可行点 \( x \)（满足 \( f_i(x) < 0 \) 和 \( Ax = b \)），设置 \( t = t^0 > 0 \)，\( \mu > 1 \)，\( \epsilon > 0 \)。

2) 中心点步骤：求解 $x^*(t) = \arg\min\ \ t f_0(x) + \phi(x) \quad \text{s.t.} \quad Ax = b$  

3) 停止准则：如果 \( \dfrac{m}{t} \leq \epsilon \)，退出；否则，令 \( x = x^*(t) \)，\( t = \mu t \)，重复。  

障碍方法通过逐渐增大 \( t \) 来逼近最优解，每个子问题需精确求解，计算成本较高。\\
2.原对偶内点法

原对偶方法同时更新原变量和对偶变量，直接求解修改的KKT条件，通常更高效。\\
2.1.修改的KKT条件

对于参数 \( t \)，中心路径点满足：$
\begin{cases}
\nabla f_0(x) + \sum_{i=1}^m \lambda_i \nabla f_i(x) + A^\top v = 0 \\
-\lambda_i f_i(x) = \dfrac{1}{t}, \quad i=1,\dots,m \\
Ax = b
\end{cases}$\\
2.2.定义残差  
\begin{align*}
r_{\text{dual}} &= \nabla f_0(x) + Df(x)^\top \lambda + A^\top v, \\
r_{\text{cent}} &= -\mathop{\text{diag}}(\lambda) f(x) - \dfrac{1}{t} \mathbf{1}, \\
r_{\text{prim}} &= Ax - b
\end{align*}

其中 \( Df(x) \) 是 \( f(x) = [f_1(x), \dots, f_m(x)]^\top \) 的雅可比矩阵。\\ 
2.3.搜索方向

原对偶搜索方向 \( \Delta y_{\text{pd}} = (\Delta x_{\text{pd}}, \Delta \lambda_{\text{pd}}, \Delta v_{\text{pd}}) \) 通过求解线性系统得到：  
    \[
    \begin{bmatrix}
    \nabla^2 f_0(x) + \sum_i \lambda_i \nabla^2 f_i(x) & Df(x)^\top & A^\top \\
    -\mathop{\text{diag}}(\lambda) Df(x) & -\mathop{\text{diag}}(f(x)) & 0 \\
    A & 0 & 0
    \end{bmatrix}
    \begin{bmatrix}
    \Delta x \\
    \Delta \lambda \\
    \Delta v
    \end{bmatrix}
    =
    -\begin{bmatrix}
    r_{\text{dual}} \\
    r_{\text{cent}} \\
    r_{\text{prim}}
    \end{bmatrix}
    \]  
2.4.代理对偶间隙

定义 \( \eta(x, \lambda) = -f(x)^\top \lambda \)。若 \( x \) 原可行且 \( \lambda, v \) 对偶可行，则 \( \eta(x, \lambda) \) 为对偶间隙。\\  
2.5.算法流程  

1) 初始化：找到 \( x \) 满足 \( f_i(x) < 0 \)（\( i=1,\dots,m \)），\( \lambda > 0 \)，设置 \( \mu > 1 \)，\( \epsilon_{\text{feas}} > 0 \)，\( \epsilon > 0 \)。  
  
2) 确定步长参数：\( t = \dfrac{\mu \cdot m}{\eta(x, \lambda)} \);

\quad 计算原对偶搜索方向 \( \Delta y_{\text{pd}} \);

\quad 直线搜索确定步长 \( s \)，更新变量组 \( y = (x, \lambda, v) \leftarrow y + s \cdot \Delta y_{\text{pd}} \)。
  
3) 停止准则：当 \( \| r_{\text{prim}} \|_2 \leq \epsilon_{\text{feas}} \)、\( \| r_{\text{dual}} \|_2 \leq \epsilon_{\text{feas}} \)，且 \( \eta(x, \lambda) \leq \epsilon \) 时，退出；否则返回步骤2。\\
原对偶方法无需精确求解子问题，迭代点未必可行，但通过牛顿法快速收敛（通常呈超线性收敛）。\\
3.寻找严格可行初始点(不可行初始点牛顿方法)

上面使用的障碍方法和原对偶内点法需要一个满足 \( f_i(x) < 0 \) 且 \( Ax = b \) 的严格可行初始点。若此类点未知，需通过如下方式寻找。  

情况1：已知 \( x^0 \in \bigcap_{i=0}^m \mathop{\text{dom}} f_i \) 满足 \( A x^0 = b \)，但可能不满足 \( f_i(x) < 0 \)。  
求解辅助问题：  
\[
\begin{cases}
\min_{s, \, x} \quad s \\
\text{s.t.} \quad f_i(x) \leq s \quad (i=1,2,\dots,m), \\
\quad\quad\,\, A x = b
\end{cases}
\]  
若最优值 \( s^* < 0 \)，则对应的 \( x^* \) 是严格可行点（因 \( f_i(x^*) \leq s^* < 0 \)）。

情况2：已知 \( x^0 \in \bigcap_{i=0}^m \mathop{\text{dom}} f_i \)，但不满足 \( A x^0 = b \)。  
  求解辅助问题：  
  \[
  \begin{cases}
  \min_{s, \, x} \quad f_0(x) \\
  \text{s.t.} \quad f_i(x) \leq s \quad (i=1,2,\dots,m), \\
  \quad\quad\,\, A x = b, \quad s = 0
  \end{cases}
  \]  
  \textbf{使用障碍方法求解}（取初始参数 \( t^0 > 0 \)）：  
  \[
  \min_{s, \, x} \quad t^0 f_0(x) - \sum_{i=1}^m \log\bigl(s - f_i(x)\bigr) 
  \quad \text{s.t.} \quad A x = b, \,\, s = 0
  \]   

情况3：不能在函数公共定义域 \( D = \bigcap_{i=0}^m \mathop{\text{dom}} f_i \) 中确定一点。  
  引入辅助变量 \( z_0, z_1, \dots, z_m \)，求解辅助问题：  
  \[
  \begin{cases}
  \min_{\substack{s, \, x, \, z_0, \dots, z_m}} \quad f_0(x + z_0) \\
  \text{s.t.} \quad f_i(x + z_i) \leq s \quad (i=1,2,\dots,m), \\
  \quad\quad\,\, A x = b, \quad s = 0, \quad z_0 = 0, \,\,\dots,\,\, z_m = 0
  \end{cases}
  \]  
  \textbf{使用障碍方法求解}（取初始参数 \( t^0 > 0 \)）：  
  \[
  \min_{\substack{s, \, x, \, z_0, \dots, z_m}} \quad t^0 f_0(x + z_0) - \sum_{i=1}^m \log\bigl(s - f_i(x + z_i)\bigr) 
  \quad \text{s.t.} \quad A x = b, \,\, s = 0, \,\, z_i = 0 \,\, (i=0,1,\dots,m)
  \]  \\
相互联系

1) 两者都是内点法，依赖于对数障碍函数和中心路径概念,中心路径是连接严格可行点与最优点的路径，参数t控制逼近精度。

2) 如果初始点不是严格可行的，两者都需要寻找严格可行初始点，求解过程中通过引入松弛变量或辅助变量确保可行性。

3）障碍方法通过求解一系列优化问题沿中心路径移动，对每个t精确求解子问题，对偶变量从中心点导出；原对偶方法直接求解修改的KKT条件，同时更新原变量和对偶变量，搜索方向由牛顿法确定，更高效且无需子问题精确解。

4) 在障碍方法中，对偶间隙为 \( \dfrac{m}{t} \)；在原对偶方法中，代理对偶间隙 \( \eta(x, \lambda) \) 用于控制停止准则。
\end{solution_cn}
% -----------------------------------------------
% 以下内容无需编辑
% -----------------------------------------------

\end{document}